% ==================================================================
% id:       lim-003
% topic:    גבולות ורציפות של פונקציות
% subtopic: משפט היינה
% solution: full
% ==================================================================
\textbf{
בכל אחד מהסעיפים הבאים, השתמשו במשפט היינה וקבעו האם קיים הגבול. \\
נמקו את תשובתכם/ן:}

\begin{enumerate}\bMs{2}
\item ${\displaystyle\lim_{x\to5}}(2x^2-3x+5)$
\ifsol
\begin{solbox}
נראה כי הגבול קיים. נשתמש במשפט היינה. לשם כך נקבע סדרה
$(a_n)_{n=1}^\infty$
כך ש-
$\displaystyle\lim_{n\to\infty}a_n=5$.
מתקיים
\begin{align*}
    \lim_{n\to\infty}f(a_n)&=\ \lim_{n\to\infty}\left(2a_n^2-3a_n+5\right)\\
    &\overset{(\ast)}{=}\ 2\cdot5^2-3\cdot5+5\\
    &=\ 40,
\end{align*}
נסיק כי הגבול המבוקש קיים ושווה ל-
$40$.

$(\ast)$ שימו לב, כי יש כאן שימוש באריתמטיקה של גבולות של סדרות.
\end{solbox}
\fi

\item $\ds{\lim_{x \to 0} e^{\frac{1}{x}}}$
\ifsol
\hint{בדקו לאן שואפת $f(a_n)$ עבור סדרות $a_n\to0$ מימין ומשמאל, והשוו.}
\begin{solbox}
נראה כי הגבול לא קיים.
תהיינה
$(a_n)_{n=1}^\infty$
ו-
$(b_n)_{n=1}^\infty$
שתי סדרות כך ש-
$$,\lim_{n\to\infty}a_n=\lim_{n\to\infty}b_n=0$$
ו-
$a_n>0$, $b_n<0$
לכל
$n\in\mathbb{N}$.
מתקיים
$$,\displaystyle\lim_{n\to\infty}f(a_n)=\lim_{n\to\infty}e^{\frac{1}{a_n}}
=\infty$$
מאחר והגבול במעריך הוא אינסוף.
מצד שני
$$,\displaystyle\lim_{n\to\infty}f(b_n)=\lim_{n\to\infty}e^{\frac{1}{b_n}}=0$$
מאחר והגבול במעריך הוא מינוס אינסוף.
התקבלו שתי תוצאות שונות, משמעו אין גבול.
\end{solbox}
\fi

\item ${\displaystyle\lim_{x\to1}}(x-\lfloor x\rfloor)$
\ifsol
\hint{בדקו מהו $\lfloor x\rfloor$ לאורך סדרות ששואפות ל-$1$ מימין ומשמאל, והשוו.}
\begin{solbox}
נפריך את קיום הגבול ע"פ ההגדרה של גבול של היינה.
נבחר שתי סדרות - הראשונה שמתכנסת ל-
$1$
בעליה, והשניה - בירידה.
לדוגמא,
$a_n=1-\tfrac{1}{n}$
ו-
$b_n=1+\tfrac{1}{n}$.
מתקיים
\begin{align*}
\displaystyle\lim_{n\to\infty}f(a_n)
&=\ \displaystyle\lim_{n\to\infty}(a_n-\lfloor a_n\rfloor)\\
&=\ \displaystyle\lim_{n\to\infty}(1-\tfrac{1}{n}-\lfloor 1-\tfrac{1}{n}\rfloor)\\
&=\ 1-0\\
&=\ 1,
\end{align*}
מצד שני
\begin{align*}
\displaystyle\lim_{n\to\infty}f(b_n)
&=\ \displaystyle\lim_{n\to\infty}(b_n-\lfloor b_n\rfloor)\\
&=\ \displaystyle\lim_{n\to\infty}(1+\tfrac{1}{n}-\lfloor 1+\tfrac{1}{n}\rfloor)\\
&=\ 1-1\\
&=\ 0,
\end{align*}
התקבלו שני גבולות שונים, לכן הגבול המבוקש לא קיים.
\end{solbox}
\fi

\item $\ds{\lim_{x \to 0} \arctan\lr{ \frac{1}{x} } }$
\ifsol
\hint{מה קורה ל-$\frac1x$ כאשר מתקרבים ל-$0$ מכל צד, ומה ידוע על $\arctan$ בפלוס ומינוס אינסוף?}
\begin{solbox}
שוב נראה כי הגבול לא קיים. נבחר שתי סדרות - הראשונה שמתכנסת ל-
$0$
מלמעלה, והשניה - מלמטה.
לדוגמא,
$a_n=-\tfrac{1}{n}$
ו-
$b_n=\tfrac{1}{n}$.
מתקיים
\begin{align*}
\displaystyle\lim_{n\to\infty}f(a_n)
&=\ \displaystyle\lim_{n\to\infty}\arctan(-n)\\
&=\ -\frac{\pi}{2},
\end{align*}
מצד שני
\begin{align*}
\displaystyle\lim_{n\to\infty}f(b_n)
&=\ \displaystyle\lim_{n\to\infty}\arctan(n)\\
&=\ \frac{\pi}{2},
\end{align*}
התקבלו שני גבולות שונים, לכן הגבול המבוקש לא קיים
(ע"פ היינה).
\end{solbox}
\fi

\item $\ds{\lim_{x \to 0} \cos\lr{\frac{1}{x} }}$
\ifsol
\hint{בחרו סדרות $a_n\to0$ שעבורן $\frac1{a_n}$ הן נקודות שבהן ערך הקוסינוס ידוע במדויק, והשוו את הגבולות.}
\begin{solbox}
נפריך את קיום הגבול בעזרת משפט היינה.
יש חששות רבים לגבי התכנסותו
של
$\cos\left(\frac{1}{x}\right)$
אשר נובעים מכך שבאינסוף אין לקוסינוס התנהגות "יציבה".

נבחר שתי סדרות ששואפות ל-
$0$.
כאן, בניגוד לסעיפים הקודמים, אין לנו יותר מדי חופש בבחירת סדרות שתפרכנה את קיום הגבול, אנו חייבים להתייחס ל"קפריזות" של קוסינוס, כלומר, עלינו לבחור סדרות בחוכמה:

שתי הסדרות חייבות לשאוף ל-
$0$,
ולאחר הרכבת הקוסינוס על הסדרות האלה כדאי לנו לדעת מה הם הערכים שיתקבלו
(הרי קוסינוס היא פונקציה "מסובכת" ואין לנו מידע לגבי ערכיה המדויקים פרט לכמה נקודות בודדות).
יודעים כי
$\cos(x)=1$
אמ"מ
$2\pi n$, $n\in\mathbb{Z}$,
ו-
$\cos(x)=0$
אמ"מ
$x=\tfrac{\pi}{2}+2\pi n$, $n\in\mathbb{Z}$.
לכן נבחר
$$a_n=\frac{1}{\tfrac{\pi}{2}+2\pi n},\quad b_n=\frac{1}{2\pi n}$$
בבירור, שתיהן שואפות ל-
$0$.
מתקיים
\begin{align*}
f(a_n)
=\cos\left(\frac{1}{a_n}\right)
= \cos\left(\tfrac{\pi}{2}+2\pi n\right)
= 0,
\end{align*}
ולכן
$\ds{\lim_{n\to\infty}}f(a_n)=0$.
מצד שני
\begin{align*}
f(b_n)
= \cos\left(\frac{1}{b_n}\right)
=\cos\left(2\pi n\right)
= 1,
\end{align*}
ולכן
$\ds{\lim_{n\to\infty}}f(b_n)=1$.
התקבלו שני גבולות שונים, לכן הגבול המבוקש לא קיים.
\end{solbox}
\fi

\item $\ds{\lim_{x \to 0} \lr{ \frac{x^2+4}{x+2} } \cdot \sin\lr{ \frac{1}{x^2} } }$
\ifsol
\hint{טפלו בכל גורם בנפרד: חשבו את הגבול של הגורם הראשון, ועבור $\sin\frac1{x^2}$ בחרו סדרות שבהן ערך הסינוס ידוע במדויק.}
\begin{solbox}
נשים לב תחילה כי
$\displaystyle\lim_{x\to0}\frac{x^2+4}{x+2}=2$,
אכן, תהא
$(c_n)_{n=1}^\infty$
סדרה ששואפת ל-
$0$.
מתקיים
\begin{align*}
    \lim_{n\to\infty}f(c_n)=\lim_{n\to\infty}\frac{c_n^2+4}{c_n+2}
    \overset{(\ast)}{=} \frac{0+4}{0+2}
    &=\ 2,
\end{align*}
כאשר המעבר
$(\ast)$
מוצדק ע"י אריתמטיקה של גבולות של סדרות.
בדומה לסעיף הקודם, נבחר
$$.a_n=\frac{1}{\sqrt{\tfrac{\pi}{2}+2\pi n}}
,\quad b_n=\frac{1}{\sqrt{\pi+2\pi n}}$$
בבירור, שתיהן שואפות ל-
$0$.
מתקיים
\begin{align*}
f(a_n)
&=\ \frac{a_n^2+4}{a_n+2}  \cdot \sin\lr{ \frac{1}{a_n^2} }\\
&=\ \frac{a_n^2+4}{a_n+2} \cdot \sin\lr{\tfrac{\pi}{2}+2\pi n}
= \frac{a_n^2+4}{a_n+2},
\end{align*}
ולכן
$\ds{\lim_{n\to\infty}}f(a_n)=2$.
מצד שני
\begin{align*}
f(b_n)
&=\ \frac{b_n^2+4}{b_n+2} \cdot \sin\lr{ \frac{1}{b_n^2} }\\
&=\ \frac{b_n^2+4}{b_n+2}  \cdot \sin\lr{\pi+2\pi n}
= 0,
\end{align*}
ולכן
$\ds{\lim_{n\to\infty}}f(b_n)=0$.
התקבלו שני גבולות שונים, לכן הגבול המבוקש לא קיים
())ע"פ היינה).
\end{solbox}
\fi
\eMs
\end{enumerate}
